\hypertarget{preserving-causal-meaning-in-a-python-identification-adapter}{%
\section{Preserving causal meaning in a Python identification
adapter}\label{preserving-causal-meaning-in-a-python-identification-adapter}}

\label{exec13:python-identification}

A common causal interface is useful only if a question keeps its meaning
when it passes into another library. An outreach decision can become a
different question if a transport drops a condition, converts an
uncertain graph into one directed graph, or treats a symbolic
probability formula as an estimated monetary return. The proposed Python
path therefore begins with a narrow service: identify an intervention
distribution in one supplied causal graph, preserve the assumptions that
make the formula relevant, and report an explicit limit when that
service cannot answer the question.

y0 is a plausible component for this service because its inspected
Python API separates a query, a mixed graph and a symbolic
identification expression. The source considered here is commit
\texttt{cc6644d3b470a1ab2a1eb9da4f0fdb0c843f5d3a}, dated 13 April 2026;
its project metadata declares development version \texttt{0.2.12-dev}.
The acquisition is later than the literature cutoff, but the reviewed
commit predates 7 October 2026. These identifiers describe inspected
source, not a tested deployment. The proposed decision is to integrate a
restricted identification path and verify its conversions before
admitting it to a prototype. Porting an entire identification engine
would create a second implementation to maintain before the first
interface has established its scientific meaning.
\citep{exec13_y0_api,exec13_y0_metadata}

\hypertarget{a-distributional-question-before-an-estimate}{%
\subsection{A distributional question before an
estimate}\label{a-distributional-question-before-an-estimate}}

Let \(A\) denote a bank action, \(Y\) an outcome and \(Z\) a collection
of explicitly named conditions. The requested object is

\[
 p(y\mid\operatorname{do}(a),z).
\]

It is distinct from an observed conditional probability, an average
difference between two actions, and a confidence interval for an
estimated difference. For a finite outcome support, a later evaluator
might form

\[
 \mu(a,z)=\sum_y y\,p(y\mid\operatorname{do}(a),z),
 \qquad \tau(z)=\mu(a_1,z)-\mu(a_0,z).
\]

That calculation requires the original outcome coordinates,
treatment/control order and units. If \(Y\) is a repayment indicator,
\(\tau\) is a probability difference. If \(Y\) is net cash received, it
has the declared currency and horizon. Multiplying the first by a loss
avoided per repayment is an additional valuation model; it is not an
operation that the identification library can infer from a variable
label.

The first adapter is deliberately limited to one scalar outcome and
treatment, optional conditions, one causal DAG or acyclic directed mixed
graph (ADMG), and declared finite discrete domains. This makes its
numerical references and interpretation concrete. It does not establish
an absence of wider functions elsewhere in y0. A graph representing
selection, transport between populations, cross-world counterfactuals,
an adaptive longitudinal regime or unresolved graph alternatives needs a
separately reviewed path. An account panel can still describe its
sampling law in the request, but a symbolic formula supplies no
independence assumption, standard error or calibration for that panel.

The scientific context should travel with the formula: target
population, action values, outcome horizon, observation and availability
clocks, variable domains, units, sign convention, sampling unit and
assumptions. A hash of that complete context identifies the request. A
changed population, a reversed action contrast or a different
availability clock produces a different identity even when the graph
picture is unchanged. The shared library, CLI and MCP operation
contracts already distinguish transport completion from scientific
disposition; the proposed adapter specializes that distinction for
identification. \citep{exec13_y0_contract}

A condition deserves particular care. Conditioning on a pretreatment
risk category and conditioning on a response to the action are different
targets. The interface should require the caller to state which is
intended, preserve the chronology, and avoid silently relabeling a
post-action variable as an adjustment covariate. Syntax can reject a
missing clock declaration. It cannot determine whether the bank's
claimed chronology or causal model is true.

\hypertarget{graph-and-coordinate-conversion}{%
\subsection{Graph and coordinate
conversion}\label{graph-and-coordinate-conversion}}

The graph object in the inspected source contains a directed component
and a component named \texttt{undirected}. In the causal mixed-graph
interpretation, the latter carries bidirected latent links: the source's
causal-graph importer maps the input type \texttt{bidirected} into that
field, and its latent expansion creates a shared hidden parent for each
such edge. It is therefore incorrect to put an unresolved CPDAG edge
into this component merely because both are stored by an undirected
graph data structure. \citep{exec13_y0_graph}

For the admitted path, a tail--arrow pair becomes a directed edge, and
an arrow--arrow pair becomes an ADMG bidirected edge. Circles,
tail--tail edges and graph collections receive an unsupported
disposition that retains the original representation. A directed edge
and a bidirected edge may coexist on the same pair of variables; erasing
one would change the model. The directed component must be acyclic, with
all endpoints in the declared node set. These checks belong in the
adapter: the inspected graph initializer sets metadata flags and does
not itself establish acyclicity.

Every variable, including an isolated one, must be passed explicitly.
Constructing a graph solely from edges can lose an outcome or an
independent condition that has no incident edge. The adapter should
compare the node and typed-edge sets before and after conversion and use
a deterministic valid topological ordering. The purpose is
reproducibility and semantic preservation, not the selection of a
statistically preferred causal ordering.

Names are another coordinate system. External identifiers may contain
spaces, non-ASCII text or labels that collide after informal
sanitization. The proposed conversion assigns opaque internal symbols
through a bijection and keeps a sidecar containing the external
identifier, displayed label, domain order, time coordinate and units. It
never guesses temporal order from a name such as \texttt{lag1}, and it
never converts category labels into numeric doses without a declared
mapping. A graph-set request remains a graph-set request even when
unsupported. An uncomputed effect set remains unknown; it is not
represented by a convenient point estimate or an empty region.

\hypertarget{why-conditional-identification-needs-a-restricted-route}{%
\subsection{Why conditional identification needs a restricted
route}\label{why-conditional-identification-needs-a-restricted-route}}

The high-level API calls ID when its conditions field is empty and IDC
otherwise. By default it converts the specific \texttt{Unidentifiable}
exception to \texttt{None}; with \texttt{strict=True} it preserves the
exception. Other exception types are not caught by that branch. The
adapter should preserve this distinction, because a malformed graph, a
runtime failure and nonidentification are different outcomes. The
inspected ID recursion supplies node/district arguments when raising
\texttt{Unidentifiable}, but those arguments should not be presented as
an independently verified graphical certificate.
\citep{exec13_y0_api,exec13_y0_id}

Two additional source findings restrict the proposed conditional route.
First, \texttt{Query.expression} builds a conditioned distribution when
conditions exist, using an \texttt{elif} branch for interventions. When
both are present, that convenience property does not add the
intervention. This is a serialization hazard, not evidence that the
separate IDC algorithm discards its treatment field:
\texttt{Identification} retains outcomes, treatments and conditions
separately. The adapter should serialize its original query triple
independently of that convenience property. \citep{exec13_y0_utils}

Second, IDC tests whether an observation may be exchanged for an action
using a separation routine. The inspected routine takes a
directed-ancestor subgraph, connects directed co-parents, flattens the
remaining edges and removes conditioning vertices. For a DAG this is a
familiar construction. The same steps do not generally implement the
collider behavior of bidirected paths. The source-level discrepancy can
be demonstrated without relying on a package's own numerical output.
\citep{exec13_y0_idc,exec13_y0_rule2,exec13_y0_separation,exec13_y0_graph}

Consider the ADMG

\[
 Y\leftrightarrow B\leftrightarrow Z,
\]

with an unrelated action \(X\). The middle vertex \(B\) is a collider.
Conditioning on it can connect \(Y\) and \(Z\); treating it as an
ordinary undirected intermediate vertex and deleting it instead
disconnects them. In the rule-2 check for exchanging \(Z\), removal of
incoming edges to isolated \(X\) and outgoing edges from \(Z\) does not
change this path. The inspected directed co-parent rule adds no
\(Y\)--\(Z\) link because there are no directed edges. Its subsequent
deletion of \(B\) therefore predicts separation where the causal mixed
graph does not.

An exact generated law makes the consequence explicit. Let \(U,V\) be
independent fair bits, let \(Y=U\), \(Z=V\), and let \(B\) report
whether \(U=V\), correctly with probability \(.9\). The observed law of
\((Y,B,Z)\) has positive probability in all eight cells. Direct
enumeration gives

\[
 P(Y=1\mid B=1,Z=1)=\frac9{10},
 \qquad
 P(Y=1\mid B=1,\operatorname{do}(Z=1))=\frac12.
\]

An intervention on \(Z\) overwrites that measurement but does not change
\(U,V\) or the mechanism generating \(B\). Exchanging the observation
for the action would therefore change the requested probability. This is
an exact reference for the scientific meaning of the graph and a trace
of the inspected source path; it is not a report of an executed y0
failure. \citep{exec13_y0_oracles}

The initial adapter can avoid that native conditional path. It first
applies ID to the joint query \(p(y,z\mid\operatorname{do}(a))\). If it
obtains an observational functional \(N(y,z,a)\), it forms

\[
 \frac{N(y,z,a)}{\sum_{y'}N(y',z,a)},
\]

on coordinates for which the denominator is positive. This is a direct
probability identity. It neither uses the questionable separation step
nor relies on the query-expression convenience property. The route is
conservative: failure to identify the joint distribution does not prove
failure to identify the conditional target. Such a failure must return
``unsupported by the admitted conditional route,'' leaving the original
question intact. Re-enabling native IDC requires a corrected or
otherwise justified path and an actual engine comparison against
independent collider references.

\hypertarget{a-front-door-reference-with-a-banking-interpretation}{%
\subsection{A front-door reference with a banking
interpretation}\label{a-front-door-reference-with-a-banking-interpretation}}

A second exact example supplies a positive identification reference. Let
\(A\) be assistance, \(M\) completion of the assisted process and \(Y\)
successful repayment. In the generated example an unobserved \(U\)
affects both assistance and repayment. The assumed graph has
\(A\to M\to Y\) and \(A\leftrightarrow Y\), with no direct effect of
assistance on repayment outside \(M\). This is a strong exclusion
assumption. In a real service workflow an agent's advice, fees or other
treatment channels could violate it, and observed associations alone do
not establish it.

Specify the finite law by

\[
 \begin{aligned}
 P(U=1)&=\tfrac12,\\
 P(A=1\mid U=u)&=\tfrac15+\tfrac35u,\\
 P(M=1\mid A=a)&=\tfrac14+\tfrac12a,\\
 P(Y=1\mid M=m,U=u)&=\tfrac1{10}+\tfrac12m+\tfrac15u.
 \end{aligned}
\]

Independent rational enumeration first constructs the observed joint
table, hiding \(U\), and evaluates

\[
 \sum_m p(m\mid a)\sum_{a'}E(Y\mid m,a')p(a').
\]

It separately evaluates the intervention law from the specified
mechanisms. Both routes produce \(13/40=.325\) under \(a=0\) and
\(23/40=.575\) under \(a=1\), for a difference of \(1/4\). The naive
observed difference is \(37/100\). The exact calculation therefore
distinguishes a successful symbolic formula from a shortcut that merely
compares observed groups. No fitted sample, randomized qualification
study or engine runtime is involved. \citep{exec13_y0_oracles}

If a proposal scenario assigns an assumed USD 1,000 benefit to each
additional successful repayment, the gross expected difference is USD
250 per account, before assistance costs and any other consequences.
That number is planning arithmetic under the declared law and value
assumption. It supplies neither observed bank ROI nor evidence that the
front-door graph describes a bank.

The inner treatment variable \(a'\) reveals an important serialization
requirement. It is a dummy summation coordinate, while the outer \(a\)
is the requested intervention value. A global string replacement could
capture the dummy variable and change the formula. The proposed result
format uses a typed expression tree with lexical bindings for
marginalized variables. It preserves probability factors, products, sums
and ratios, alpha-renames bound symbols when needed, and maps free
coordinates back to the original question. The inspected symbolic
classes expose those structures; rendered LaTeX or the library's
Python-like display string is suitable for presentation, not as an
executable interchange format. \citep{exec13_y0_dsl}

An independent negative reference is equally necessary. In a graph
allowing \(A\to Y\) and \(A\leftrightarrow Y\), take a fair hidden
\(U\), let \(A=U\), and let an independent bit \(E\) equal one with
probability \(.2\). The mechanisms \(Y=A\mathbin{\mathrm{xor}}E\) and
\(Y=U\mathbin{\mathrm{xor}}E\) give the same strictly positive observed
law of \((A,Y)\). Under intervention their effects are respectively
\(.6\) and zero. Both belong to the model allowing those edges; no
faithfulness restriction is imposed. A library cannot choose between the
effects from their identical observed distribution. A successful adapter
retains this nonidentification rather than translating every completed
request into a finite effect. \citep{exec13_y0_oracles}

\hypertarget{responses-deployment-and-the-next-prototype-decision}{%
\subsection{Responses, deployment and the next prototype
decision}\label{responses-deployment-and-the-next-prototype-decision}}

A symbolic identification result should contain the original query,
graph and context identities, engine revision, typed functional, support
conditions and source provenance. Its numeric estimate, standard error
and confidence region remain absent. When the unconditional ID path
reports nonidentification in an admitted model, the response should
identify that origin and retain its diagnostic without exaggerating its
verification. When the conditional joint route fails, the response is
unsupported. An invalid graph is rejected; a missing engine, internal
failure, timeout or cancellation is unavailable. These distinctions
determine whether the next action is to repair an input, add scientific
information, select a different method or provision a runtime.

Library, CLI and MCP surfaces should call the same core implementation
and return the same context and result identities. The symbolic service
requires graph and question artifacts, not bank account rows. A worker
can enforce wall-time, memory and output budgets while keeping engine
diagnostics separate from scientific conclusions. Cancellation must
preserve a completed result if publication has already occurred; an
interrupted partial expression cannot become an identified functional.
These are proposed implementation requirements, not performance
guarantees supplied by the inspected y0 API.

The pinned project metadata requires Python 3.11 or later and lists
twelve direct runtime dependencies, including NetworkX, SciPy, pandas,
scikit-learn and statsmodels. The build backend is \texttt{uv\_build} in
its declared version range. Its optional R extra uses \texttt{rpy2};
upstream test groups include additional packages that are not all
runtime requirements of the proposed pure-Python path. The inspected
tree has no root lockfile establishing resolved transitive versions. A
prototype admission therefore needs a frozen interpreter/platform and
dependency manifest, saved distribution hashes, actual build/import logs
and license notices. Source inspection alone cannot promise
compatibility with whatever environment happens to be installed.
\citep{exec13_y0_metadata,exec13_y0_tree}

The saved BSD three-clause license permits redistribution under its
notice, disclaimer and nonendorsement conditions. That supports a
possible integration route, while leaving dependency rights and actual
deployment behavior to be checked. It supplies no scientific warranty. A
port may become justified if a later measurement shows that the Python
environment or latency is a decisive constraint; there is no such
measured advantage for porting in the present evidence.
\citep{exec13_y0_license}

The prototype's first engine trials should compare the actual pinned
runtime with the independent front-door, bow and collider references,
including retained unsupported requests. Conversion checks should
exercise isolated variables, edge-type distinctions, coordinate
permutations, dummy-variable binding, zero conditional support and
distinct failures. A successful trial would support the named adapter in
its declared regime. It would not establish discovery accuracy, graph
truth, estimation calibration or institutional acceptance. The value of
this restricted design is that each later claim has a specific test and
a specific scientific object that must survive the interface.
\label{exec13:python-identification-end}
